\clearpage
\section{The Propers: Themes and Movement}
\zlabel{triptych:concise:themes:start}
The King who stands before Pilate is the Son through whom every creature exists. Colossians identifies his death as the means of reconciliation. The Mass celebrates his redemption of those who become his people: they receive grace to hear him, peace with God and one another, and the hope of sharing his enduring reign.

\subsection{The slain Lamb and the divided family}
The Introit acclaims the Lamb as slain. In Apocalypse 5 he is worthy to open the sealed book because he was slain and redeemed a people. The heavenly worshippers praise him in the words the chant selects. Apocalypse 1 likewise names kingdom and priesthood in its account of his faithful witness and saving blood. The assembly praises this sacrificial death before hearing the Gospel's account of the trial that precedes it.

The Collect asks for the restoration of all things in the beloved Son, then names the human families scattered by sin. Their division concerns their relation to Christ and their relations with one another. The prayer asks that they be gathered under his dominion, described as gentle. The Introit's psalm verse asks God for judgment and justice for the King. Both prayers concern the rule under which these peoples are to live.

\subsection{The Creator stands before a governor}
Colossians gives thanks for deliverance into the Son's kingdom before identifying him as Creator, head of the Church and firstborn from the dead. Chrysostom distinguishes giving the inheritance from making its recipients fit to share it. God changes the recipients as well as rescuing them, and transfer into the kingdom admits them to life with the Son. Chrysostom also explains why the title firstborn cannot make him a creature: all visible and invisible things were created through him and continue to depend on him. The same Son becomes head of the Church and makes peace through his blood. Those he reconciles already depend on him for existence.\footnote{Chrysostom, \emph{Homilies on Colossians} II, on 1:12--15; III, on 1:15--20.}

The Gradual sings dominion over every nation and ruler; the Alleluia, a kingdom that endures. Daniel's human figure receives dominion before the Ancient of Days. Jerome explains this gift according to the humanity the Son assumed and the self-emptying described in Philippians 2. Received power does not imply an inferior divine nature: the eternal Son truly becomes man and receives the kingdom in that condition. The saints share his heavenly kingdom. Its scope and permanence exceed the earthly kingdoms that arise and pass away in the vision.

Before Pilate, Jesus says his kingdom is not of this world. Augustine distinguishes its origin from its presence: believers remain in the world but become Christ's people through regeneration. He cites the Epistle's transfer out of darkness to explain that change. At the Gospel's close, belonging to truth precedes hearing the King's voice. All people were created through the Word, yet not all believe. Augustine therefore understands the hearing named here as a further gift, inward belief and obedience made possible by grace. Hearing follows that gift; it does not earn the hearer's first belonging to truth.\footnote{Augustine, \emph{Tractates on John} 115.2,4.}

\clearpage
\subsection{Intercession for the nations}
The Offertory quotes the Father's command to the Son to ask for the nations as his inheritance. Augustine understands that asking through the humanity Christ assumed, his sacrifice and his intercession. The Son asks for the salvation of those he receives. Read in this way, the peoples for whom the Collect seeks healing are the inheritance for which their King intercedes. His petition concerns the nations to the ends of the earth.\footnote{Augustine, \emph{Expositions of the Psalms}, Ps 2.7.}

In the Secret, the Father is asked that the Son offered in sacrifice bestow unity and peace upon those nations. Christ is both the one offered and the giver of the gifts sought. The Epistle identifies the blood of his Cross as the means of peace. These passages can be understood together as the reconciliation accomplished by his death and sought for the peoples who remain divided, including those making the offering.

The Preface names the Son eternal Priest and universal King, recounts his self-offering, and describes the kingdom through truth, life, holiness, grace, justice, love and peace. The assembly then gives thanks with the angels. Chrysostom asks why heaven needs reconciliation when it is the model of obedience in the Lord's Prayer. The holy angels have not sinned; their estrangement concerns sinful humanity's dishonor of their Lord. Christ brings humanity into their company. In light of that account, the shared thanksgiving can be read as the fellowship his death restores.\footnote{Chrysostom, \emph{Homilies on Colossians} III, on reconciliation of things in heaven.}

Bellarmine distinguishes three senses of the Gradual's universal homage: widespread conversion allowing exceptions, complete acknowledgment at judgment, and Christ's present authority even over those who resist. Conversion changes a person's allegiance; acknowledgment at judgment cannot be evaded; authority already belongs to Christ while subjects resist him. His right over every nation does not mean every person already obeys willingly. The Collect and Secret continue to ask for the healing these peoples need.\footnote{Bellarmine, \emph{Commentary on Psalms}, Ps 71, numbered comment 10 on biblical v. 11.}

\subsection{Peace, service and shared reign}
The Communion sings the King's blessing of peace at the close of a psalm describing his voice amid a storm. Augustine reads its conclusion as the Church's preservation through the world's flood and future reign with Christ. Peace is received amid continuing adversity. Bellarmine likewise distinguishes strength on the present journey from the final blessing toward which it leads.\footnote{Augustine, \emph{Expositions of the Psalms}, Ps 28.10--11; Bellarmine, Ps 28, conclusion.}

Chrysostom applies his account of reconciliation to the Church's exchange of peace. The greeting is addressed to all, yet someone receives it while keeping peace with one neighbor and a feud with another. He compares this selective peace to a body with some members at rest and others disordered. Its health requires concord throughout the body. The Son came to reconcile heaven and earth, so deliberate enmity toward a fellow member contradicts the gift received. In this account, the peace exchanged in worship requires the healing of the relationship that remains divided.

The Postcommunion asks that those already nourished by the sacrament, who glory in serving under Christ's standard, may reign with him in heaven. Their service can be understood as perseverance toward that inheritance. Read beside the Gospel, the military image concerns fidelity to the King whose servants do not fight to prevent his Passion. Such service is consistent with the peace sought at the altar. The nourishment already received can sustain the people who serve him through present conflict toward the lasting reign they ask to share.
\zlabel{triptych:concise:themes:end}
